\documentclass[a4paper, 11pt]{article}

\usepackage[top=3cm, left=2cm, text={17cm, 24cm}]{geometry}

\usepackage[hidelinks]{hyperref}

\usepackage[slovak]{babel}


\begin{document}
    \begin{titlepage}
        \begin{center}
            {\Huge Vysoké učení technické v Brně} \\ 
            {\huge Fakulta informačních technologií}  \\
            
            \vspace{\stretch{0.382}}
            {\Huge Interpreter SOL26}\\ 
            {\large Dokumentácia Implementácie}

            \vspace{\stretch{0.618}}
        \end{center}

        {\Large \today \hfill Kristián Lupták}
    \end{titlepage}

    \tableofcontents
    \newpage
    

    \section{Úvod}
        Tento dokument popisuje implementáciu Interpreteru. 
        Vysvetľuje ako je~program rozdelený do~častí pričom informuje ako je~každá časť navrhnutá, aké objektovo orientované princípy používa a~ako spolupracuje s~inými časťami.
        Dokument je obohatený UML diagramom celého objektového návrhu. \\

        \noindent
        V~dokumente sa~taktiež nachádzajú doplňujúce informácie o~vývoji projektu(napr. používanie AI) a~taktiež ako by sa~teoreticky dal Interpret rozšíriť.
        \medskip
    \section{Popis návrhu Interpreteru}
        Interpreter je logicky rozdelený na viacero častí pričom každá nasleduje až po dokončení predchodzej.

        \subsection{Sémantická analýza}
        Ako prvá sa najprv uskutoční \textbf{kompletná} statická sémantická analýza, ktorá skontroluje vstupný program a pri odhalení chyby vracia korešpondujúji chybový kód. 
        Tento princíp umožňuje odchytiť niektoré chyby už na začiatku spustenia Interpreteru pred výkonávaním samotného programu.
        Pri vyhotovovaní následných už statické chyby \textbf{niesú} vôbec kontrolované. \\

        \subsection{Príprava prostredia na Interpretáciu}
        Následne sa sa vytvoria objekty typu \textbf{SolClass}(viď \ref{}) pre vstavané Triedy, a vytvoria sa ich metódy typu \texttt{SolMethod}, ktoré špecificky obsahujú odkaz na implementáciu metódy. 
        Po vytvoren9 vstavaných tried sa rekurzívne vytvoria Triedy definované používateľom pričom ich metódy sú taktiež typu \texttt{SolMethod}, ktoré ale neobsahujú odkaz na funkciu implementácie ale vkladá sa do nich blok (postupnosť priradení) metódy.
        Ďalej sa vytvoria singleton objekty \texttt{true, false, nil} a globálne rozhranie, do ktorého sú tieto objekty vložené. Nakoniec sa vytvorí objekt typu \texttt{Main}, ktorému je poslaná správa \texttt{run}.
        
        \subsection{Vyhodnocovanie Programu}
        Program sa vykonáva na princípe zasielaní správ, kde správa obsahuje základné argumenty ako \texttt{príjemca}, \texttt{selektor} a \texttt{argumenty}, ale je obohatený aj typom správy a triednym kontextom.
        Typ správy závisí od toho ako bol vyhodnotený objekt príjemcu, to znamená, že či bol objekt zavolaný cez \texttt{self}, \text{super} alebo premennou.
        Triedny kontext je trieda, ktorá definuje danú metódu, v prípade bloku je to trieda, kde bol blok \texttt{definovaný}.
        Pri invokovaní metódy skrytej za selektorom sa rozlišuje či je daná metóda vstavaná, čo spôsobí invokáciu definovanej metódy, alebo je metóda definovaná uživateľom, čo spôsobí invokáciu bloku uloženého v metóde. 
        Pri invokácii bloku program postupne prechádza jednotlivé priradenia a vyhodnocovanie ich výrazov následne spúšta.
        Jednotlivé výrazy sú následne výhodnocované na základe ich typu pričom Interpret spúšťa ich korespondné vyhodnocovacie funkcie.

        \medskip
\end{document}